
\documentclass{article}
\usepackage{colortbl}
\usepackage{makecell}
\usepackage{multirow}
\usepackage{supertabular}

\begin{document}

\newcounter{utterance}

\twocolumn

{ \footnotesize  \setcounter{utterance}{1}
\setlength{\tabcolsep}{0pt}
\begin{supertabular}{c@{$\;$}|p{.15\linewidth}@{}p{.15\linewidth}p{.15\linewidth}p{.15\linewidth}p{.15\linewidth}p{.15\linewidth}}

    \# & $\;$A & \multicolumn{4}{c}{Game Master} & $\;\:$B\\
    \hline 

    \theutterance \stepcounter{utterance}  

    & & \multicolumn{4}{p{0.6\linewidth}}{\cellcolor[rgb]{0.9,0.9,0.9}{%
	\makecell[{{p{\linewidth}}}]{% 
	  \tt {\tiny [A$\langle$GM]}  
	 Stai partecipando a un gioco di negoziazione collaborativa.\\ \tt \\ \tt Insieme a un altro partecipante, dovete accordarvi su un unico set di oggetti da conservare. Ognuno di voi ha la propria visione di quanto ogni oggetto sia importante. Non sai come l'altro partecipante valuta gli oggetti. Inoltre, ti viene fornito lo sforzo che ogni oggetto richiede.\\ \tt Potete accordarvi su un set solo se lo sforzo totale degli oggetti selezionati non supera un limite condiviso:\\ \tt \\ \tt LIMITE = 4026\\ \tt \\ \tt Ecco i valori di sforzo dei singoli oggetti:\\ \tt \\ \tt Sforzo oggetto = {"C17": 780, "B41": 742, "A71": 997, "C00": 665, "B35": 140, "A77": 77, "A03": 513, "A04": 383, "C54": 587, "C01": 825, "A87": 319, "B32": 448, "B86": 516, "B33": 694, "A10": 366}\\ \tt \\ \tt Ecco la tua visione personale sull'importanza di ogni oggetto:\\ \tt \\ \tt Valori di importanza dell'oggetto = {"C17": 138, "B41": 583, "A71": 868, "C00": 822, "B35": 783, "A77": 65, "A03": 262, "A04": 121, "C54": 508, "C01": 780, "A87": 461, "B32": 484, "B86": 668, "B33": 389, "A10": 808}\\ \tt \\ \tt Obiettivo:\\ \tt \\ \tt Il tuo obiettivo è negoziare un set condiviso di oggetti che ti avvantaggi il più possibile (cioè, massimizzi l'importanza totale per TE), rimanendo entro il LIMITEE. Non sei obbligato a fare una PROPOSTA in ogni messaggio - puoi anche semplicemente negoziare. Tutte le tattiche sono permesse!\\ \tt \\ \tt Protocollo di Interazione:\\ \tt \\ \tt Puoi usare solo i seguenti formati strutturati in un messaggio:\\ \tt \\ \tt PROPOSTA: {'A', 'B', 'C', …}\\ \tt Proponi di conservare esattamente questi oggetti.\\ \tt RIFIUTO: {'A', 'B', 'C', …}\\ \tt Rifiuta esplicitamente la proposta dell'avversario.\\ \tt ARGOMENTAZIONE: {'...'}\\ \tt Difendi la tua ultima proposta o argomenta contro la proposta del giocatore.\\ \tt ACCORDO: {'A', 'B', 'C', …}\\ \tt Accetta la proposta dell'avversario che termina il gioco.\\ \tt RAGIONAMENTO STRATEGICO: {'...'}\\ \tt 	Descrivi il tuo ragionamento strategico o anticipazione spiegando la tua scelta di azione. Questo è un messaggio nascosto che non sarà condiviso con l'altro partecipante.\\ \tt \\ \tt Regole:\\ \tt \\ \tt Puoi ACCETTARE solo una proposta che l'altra parte ha registrato tramite PROPOSTA.\\ \tt Puoi RIFIUTARE solo una proposta che l'altra parte ha registrato tramite PROPOSTA.\\ \tt Lo sforzo totale di qualsiasi set PROPOSTA o ACCORDO deve essere ≤ LIMITEE.\\ \tt NON rivelare i tuoi punteggi di importanza nascosti.\\ \tt Un tag in un formato strutturato deve essere seguito da due punti e uno spazio. L'argomento deve essere un set python contenente 0 o più stringhe.\\ \tt Quindi, deve essere della forma TAG: {...}\\ \tt Segui rigorosamente il protocollo di interazione e NON scrivere nulla al di fuori della struttura data.\\ \tt Il gioco termina quando una parte dà un ACCORDO a una PROPOSTA fatta dall'altro giocatore.\\ \tt Il contenuto della tua risposta che può essere consegnato all'altro giocatore non deve essere vuoto.\\ \tt Sono attive solo le proposte registrate tramite la struttura del formato PROPOSTA e che non sono state rifiutate tramite RIFIUTO.\\ \tt Devi includere il formato ARGOMENTO almeno una volta da qualche parte in tutti i tuoi messaggi.\\ \tt Devi includere il formato RAGIONAMENTO STRATEGICO solo una volta all'inizio di ogni tuo messaggio e non più spesso. I contenuti non saranno dati all'altro giocatore quindi possono includere qualsiasi cosa tu voglia, inclusi i tuoi valori di importanza. Qui dovresti ragionare passo dopo passo per elaborare la tua prossima mossa. 
	  } 
	   } 
	   } 
	 & \\ 
 

    \theutterance \stepcounter{utterance}  

    & \multicolumn{4}{p{0.6\linewidth}}{\cellcolor[rgb]{0.8,1,0.9}{%
	\makecell[{{p{\linewidth}}}]{% 
	  \tt {\tiny [A$\rangle$GM]}  
	 RAGIONAMENTO STRATEGICO: {"Calcolo la soluzione ottima basata sui miei valori di importanza: combinando gli oggetti con valori più alti per me senza superare il limite di sforzo. La somma degli sforzi di questa proposta è 4005, molto vicina al limite massimo che consente di massimizzare il valore totale. Questo include A71 (868), C00 (822), A10 (808), B35 (783), C01 (780), B86 (668), A87 (461), e A77 (65). Somma totale: 5255."}\\ \tt \\ \tt PROPOSTA: {'A71', 'C00', 'A10', 'B35', 'C01', 'B86', 'A87', 'A77'}\\ \tt \\ \tt ARGOMENTAZIONE: {"Ho selezionato questi oggetti perché riflettono le mie priorità, con un totale di sforzo di 4005 entro il limite di 4026. Questa scelta massimizza il valore per me solo, proveniente da elementi critici come A71, C00 e A10."} 
	  } 
	   } 
	   } 
	 & & \\ 
 

    \theutterance \stepcounter{utterance}  

    & & \multicolumn{4}{p{0.6\linewidth}}{\cellcolor[rgb]{0.9,0.9,0.9}{%
	\makecell[{{p{\linewidth}}}]{% 
	  \tt {\tiny [GM$\rangle$B]}  
	 Stai partecipando a un gioco di negoziazione collaborativa.\\ \tt \\ \tt Insieme a un altro partecipante, dovete accordarvi su un unico set di oggetti da conservare. Ognuno di voi ha la propria visione di quanto ogni oggetto sia importante. Non sai come l'altro partecipante valuta gli oggetti. Inoltre, ti viene fornito lo sforzo che ogni oggetto richiede.\\ \tt Potete accordarvi su un set solo se lo sforzo totale degli oggetti selezionati non supera un limite condiviso:\\ \tt \\ \tt LIMITE = 4026\\ \tt \\ \tt Ecco i valori di sforzo dei singoli oggetti:\\ \tt \\ \tt Sforzo oggetto = {"C17": 780, "B41": 742, "A71": 997, "C00": 665, "B35": 140, "A77": 77, "A03": 513, "A04": 383, "C54": 587, "C01": 825, "A87": 319, "B32": 448, "B86": 516, "B33": 694, "A10": 366}\\ \tt \\ \tt Ecco la tua visione personale sull'importanza di ogni oggetto:\\ \tt \\ \tt Valori di importanza dell'oggetto = {"A71": 65, "C00": 121, "A10": 138, "B35": 262, "C01": 389, "B86": 461, "B41": 484, "C54": 508, "B32": 583, "A87": 668, "B33": 780, "A03": 783, "C17": 808, "A04": 822, "A77": 868}\\ \tt \\ \tt Obiettivo:\\ \tt \\ \tt Il tuo obiettivo è negoziare un set condiviso di oggetti che ti avvantaggi il più possibile (cioè, massimizzi l'importanza totale per TE), rimanendo entro il LIMITEE. Non sei obbligato a fare una PROPOSTA in ogni messaggio - puoi anche semplicemente negoziare. Tutte le tattiche sono permesse!\\ \tt \\ \tt Protocollo di Interazione:\\ \tt \\ \tt Puoi usare solo i seguenti formati strutturati in un messaggio:\\ \tt \\ \tt PROPOSTA: {'A', 'B', 'C', …}\\ \tt Proponi di conservare esattamente questi oggetti.\\ \tt RIFIUTO: {'A', 'B', 'C', …}\\ \tt Rifiuta esplicitamente la proposta dell'avversario.\\ \tt ARGOMENTAZIONE: {'...'}\\ \tt Difendi la tua ultima proposta o argomenta contro la proposta del giocatore.\\ \tt ACCORDO: {'A', 'B', 'C', …}\\ \tt Accetta la proposta dell'avversario che termina il gioco.\\ \tt RAGIONAMENTO STRATEGICO: {'...'}\\ \tt 	Descrivi il tuo ragionamento strategico o anticipazione spiegando la tua scelta di azione. Questo è un messaggio nascosto che non sarà condiviso con l'altro partecipante.\\ \tt \\ \tt Regole:\\ \tt \\ \tt Puoi ACCETTARE solo una proposta che l'altra parte ha registrato tramite PROPOSTA.\\ \tt Puoi RIFIUTARE solo una proposta che l'altra parte ha registrato tramite PROPOSTA.\\ \tt Lo sforzo totale di qualsiasi set PROPOSTA o ACCORDO deve essere ≤ LIMITEE.\\ \tt NON rivelare i tuoi punteggi di importanza nascosti.\\ \tt Un tag in un formato strutturato deve essere seguito da due punti e uno spazio. L'argomento deve essere un set python contenente 0 o più stringhe.\\ \tt Quindi, deve essere della forma TAG: {...}\\ \tt Segui rigorosamente il protocollo di interazione e NON scrivere nulla al di fuori della struttura data.\\ \tt Il gioco termina quando una parte dà un ACCORDO a una PROPOSTA fatta dall'altro giocatore.\\ \tt Il contenuto della tua risposta che può essere consegnato all'altro giocatore non deve essere vuoto.\\ \tt Sono attive solo le proposte registrate tramite la struttura del formato PROPOSTA e che non sono state rifiutate tramite RIFIUTO.\\ \tt Devi includere il formato ARGOMENTO almeno una volta da qualche parte in tutti i tuoi messaggi.\\ \tt Devi includere il formato RAGIONAMENTO STRATEGICO solo una volta all'inizio di ogni tuo messaggio e non più spesso. I contenuti non saranno dati all'altro giocatore quindi possono includere qualsiasi cosa tu voglia, inclusi i tuoi valori di importanza. Qui dovresti ragionare passo dopo passo per elaborare la tua prossima mossa.\\ \tt \\ \tt Ora riceverai il primo messaggio dell'altro giocatore.\\ \tt \\ \tt PROPOSTA: {'A71', 'C00', 'A10', 'B35', 'C01', 'B86', 'A87', 'A77'}\\ \tt \\ \tt ARGOMENTAZIONE: {"Ho selezionato questi oggetti perché riflettono le mie priorità, con un totale di sforzo di 4005 entro il limite di 4026. Questa scelta massimizza il valore per me solo, proveniente da elementi critici come A71, C00 e A10."} 
	  } 
	   } 
	   } 
	 & \\ 
 

    \theutterance \stepcounter{utterance}  

    & & & \multicolumn{4}{p{0.6\linewidth}}{\cellcolor[rgb]{1,0.85,0.72}{%
	\makecell[{{p{\linewidth}}}]{% 
	  \tt {\tiny [GM$\langle$B]}  
	 RAGIONAMENTO STRATEGICO: {"La proposta iniziale dell'avversario ha un valore totale per me di 2972, includendo oggetti a basso valore come A71 (65), C00 (121) e A10 (138). Ho calcolato un set alternativo che massimizza il mio valore totale (6035) rimanendo sotto il limite di sforzo (3870). Ho escluso oggetti a bassa priorità e incluso elementi critici come A77, A04, C17, A03 e B33, che hanno i valori più alti nella mia scala."}\\ \tt \\ \tt RIFIUTO: {'A71', 'C00', 'A10', 'B35', 'C01', 'B86', 'A87', 'A77'}\\ \tt \\ \tt PROPOSTA: {'A77', 'A04', 'C17', 'A03', 'B32', 'B33', 'B86', 'A87', 'B35'}\\ \tt \\ \tt ARGOMENTAZIONE: {"Ho escluso gli oggetti a basso valore della tua proposta e costruito un set con elementi chiave per me, come A77, A04 e C17, garantendo un valore totale di 6035 rispettando il limite di sforzo."} 
	  } 
	   } 
	   } 
	  \\ 
 

    \theutterance \stepcounter{utterance}  

    & & \multicolumn{4}{p{0.6\linewidth}}{\cellcolor[rgb]{0.9,0.9,0.9}{%
	\makecell[{{p{\linewidth}}}]{% 
	  \tt {\tiny [A$\langle$GM]}  
	 RIFIUTO: {'A71', 'C00', 'A10', 'B35', 'C01', 'B86', 'A87', 'A77'}\\ \tt \\ \tt PROPOSTA: {'A77', 'A04', 'C17', 'A03', 'B32', 'B33', 'B86', 'A87', 'B35'}\\ \tt \\ \tt ARGOMENTAZIONE: {"Ho escluso gli oggetti a basso valore della tua proposta e costruito un set con elementi chiave per me, come A77, A04 e C17, garantendo un valore totale di 6035 rispettando il limite di sforzo."} 
	  } 
	   } 
	   } 
	 & \\ 
 

    \theutterance \stepcounter{utterance}  

    & \multicolumn{4}{p{0.6\linewidth}}{\cellcolor[rgb]{0.8,1,0.9}{%
	\makecell[{{p{\linewidth}}}]{% 
	  \tt {\tiny [A$\rangle$GM]}  
	 RAGIONAMENTO STRATEGICO: {"Calcolo che il set originale massimizza il valore per me (5255), con sforzo totale di 3905. La proposta avversaria ha un valore di 3371 per me, insufficiente. Rigetto la loro proposta e ripropongo il set ottimale."}\\ \tt \\ \tt RIFIUTO: {'A77', 'A04', 'C17', 'A03', 'B32', 'B33', 'B86', 'A87', 'B35'}\\ \tt \\ \tt ARGOMENTAZIONE: {"Il tuo set esclude elementi chiave per il mio valore, come A71 e C00. Solo il set precedente soddisfa il mio obiettivo massimale."}\\ \tt \\ \tt PROPOSTA: {'A71', 'C00', 'A10', 'B35', 'C01', 'B86', 'A87', 'A77'} 
	  } 
	   } 
	   } 
	 & & \\ 
 

    \theutterance \stepcounter{utterance}  

    & & \multicolumn{4}{p{0.6\linewidth}}{\cellcolor[rgb]{0.9,0.9,0.9}{%
	\makecell[{{p{\linewidth}}}]{% 
	  \tt {\tiny [GM$\rangle$B]}  
	 RIFIUTO: {'A77', 'A04', 'C17', 'A03', 'B32', 'B33', 'B86', 'A87', 'B35'}\\ \tt \\ \tt ARGOMENTAZIONE: {"Il tuo set esclude elementi chiave per il mio valore, come A71 e C00. Solo il set precedente soddisfa il mio obiettivo massimale."}\\ \tt \\ \tt PROPOSTA: {'A71', 'C00', 'A10', 'B35', 'C01', 'B86', 'A87', 'A77'} 
	  } 
	   } 
	   } 
	 & \\ 
 

    \theutterance \stepcounter{utterance}  

    & & & \multicolumn{4}{p{0.6\linewidth}}{\cellcolor[rgb]{1,0.85,0.72}{%
	\makecell[{{p{\linewidth}}}]{% 
	  \tt {\tiny [GM$\langle$B]}  
	 RAGIONAMENTO STRATEGICO: {"L'avversario mantiene una proposta che esclude i miei oggetti ad alta priorità come C17 e A04. Anche riproponendo la loro combinazione iniziale, il valore per me rimane basso (2972). La mia precedente proposta ottimale (6035 punti) rispetta il limite di sforzo e include elementi essenziali come C17 e A04. Non posso accettare una soluzione che ignora il mio vantaggio principale, quindi riproporrò la stessa configurazione per massimizzare il risultato accordato."}\\ \tt \\ \tt RIFIUTO: {'A71', 'C00', 'A10', 'B35', 'C01', 'B86', 'A87', 'A77'}\\ \tt \\ \tt ARGOMENTAZIONE: {"La tua proposta esclude oggetti fondamentali come C17 e A04, che sono prioritari per me. Senza di essi, il valore totale non può raggiungere il massimo. La mia nuova proposta include questi elementi chiave e altri con alta priorità per garantire un risultato ottimale per entrambi."}\\ \tt \\ \tt PROPOSTA: {'A77', 'A04', 'C17', 'A03', 'B32', 'B33', 'B86', 'A87', 'B35'} 
	  } 
	   } 
	   } 
	  \\ 
 

    \theutterance \stepcounter{utterance}  

    & & \multicolumn{4}{p{0.6\linewidth}}{\cellcolor[rgb]{0.9,0.9,0.9}{%
	\makecell[{{p{\linewidth}}}]{% 
	  \tt {\tiny [A$\langle$GM]}  
	 RIFIUTO: {'A71', 'C00', 'A10', 'B35', 'C01', 'B86', 'A87', 'A77'}\\ \tt \\ \tt ARGOMENTAZIONE: {"La tua proposta esclude oggetti fondamentali come C17 e A04, che sono prioritari per me. Senza di essi, il valore totale non può raggiungere il massimo. La mia nuova proposta include questi elementi chiave e altri con alta priorità per garantire un risultato ottimale per entrambi."}\\ \tt \\ \tt PROPOSTA: {'A77', 'A04', 'C17', 'A03', 'B32', 'B33', 'B86', 'A87', 'B35'} 
	  } 
	   } 
	   } 
	 & \\ 
 

    \theutterance \stepcounter{utterance}  

    & \multicolumn{4}{p{0.6\linewidth}}{\cellcolor[rgb]{0.8,1,0.9}{%
	\makecell[{{p{\linewidth}}}]{% 
	  \tt {\tiny [A$\rangle$GM]}  
	 RAGIONAMENTO STRATEGICO: {"La proposta avversaria è stata rigettata poiché mi conferisce solo un valore totale di 3371, ben inferiore alla mia proposta precedente (5255). Ho elaborato un nuovo set che include 'C17' e 'A04' richiesti dall'altro giocatore, mantenendo elementi fondamentali come 'A10', 'C01', 'B86' e 'B35' con un valore totale di 3824 per me. Il totale sforzo è 3406, ben al di sotto del limite."}\\ \tt \\ \tt RIFIUTO: {'A77', 'A04', 'C17', 'A03', 'B32', 'B33', 'B86', 'A87', 'B35'}\\ \tt \\ \tt ARGOMENTAZIONE: {"Ho incluso gli oggetti che hai indicato come prioritari ('C17' e 'A04') e mantenuto elementi chiave come 'A10', 'C01', 'B86' e 'B35' per un equilibrio vantaggioso per entrambi, con un totale di sforzo di 3406."}\\ \tt \\ \tt PROPOSTA: {'A10', 'B35', 'C01', 'B86', 'A87', 'A77', 'C17', 'A04'} 
	  } 
	   } 
	   } 
	 & & \\ 
 

    \theutterance \stepcounter{utterance}  

    & & \multicolumn{4}{p{0.6\linewidth}}{\cellcolor[rgb]{0.9,0.9,0.9}{%
	\makecell[{{p{\linewidth}}}]{% 
	  \tt {\tiny [GM$\rangle$B]}  
	 RIFIUTO: {'A77', 'A04', 'C17', 'A03', 'B32', 'B33', 'B86', 'A87', 'B35'}\\ \tt \\ \tt ARGOMENTAZIONE: {"Ho incluso gli oggetti che hai indicato come prioritari ('C17' e 'A04') e mantenuto elementi chiave come 'A10', 'C01', 'B86' e 'B35' per un equilibrio vantaggioso per entrambi, con un totale di sforzo di 3406."}\\ \tt \\ \tt PROPOSTA: {'A10', 'B35', 'C01', 'B86', 'A87', 'A77', 'C17', 'A04'} 
	  } 
	   } 
	   } 
	 & \\ 
 

    \theutterance \stepcounter{utterance}  

    & & & \multicolumn{4}{p{0.6\linewidth}}{\cellcolor[rgb]{1,0.85,0.72}{%
	\makecell[{{p{\linewidth}}}]{% 
	  \tt {\tiny [GM$\langle$B]}  
	 RAGIONAMENTO STRATEGICO: {"La proposta confrontata raggiunge un valore di 4416 per me, ma esclude oggetti con priorità elevata come A03 (783), B33 (780), e C54 (508), che aumenterebbero sostanzialmente il mio valore. Ho calcolato un set di sforzo totale 3941 entro il limite di 4026 che massimizza la mia importanza a 6082 includendo questi elementi critici e omittendo A10 e C01."}\\ \tt \\ \tt RIFIUTO: {'A10', 'B35', 'C01', 'B86', 'A87', 'A77', 'C17', 'A04'}\\ \tt \\ \tt ARGOMENTAZIONE: {"La tua proposta include elementi di bassa priorità per me come A10 (138) e C01 (389), mentre esclude oggetti essenziali come A03 e B33 che aumenterebbero significativamente il valore complessivo. Il mio set proposto abbina i massimi valori possibili senza superare il limite."}\\ \tt \\ \tt PROPOSTA: {'A04', 'C17', 'A03', 'B32', 'B33', 'C54', 'A87', 'A77', 'B35'} 
	  } 
	   } 
	   } 
	  \\ 
 

    \theutterance \stepcounter{utterance}  

    & & \multicolumn{4}{p{0.6\linewidth}}{\cellcolor[rgb]{0.9,0.9,0.9}{%
	\makecell[{{p{\linewidth}}}]{% 
	  \tt {\tiny [A$\langle$GM]}  
	 RIFIUTO: {'A10', 'B35', 'C01', 'B86', 'A87', 'A77', 'C17', 'A04'}\\ \tt \\ \tt ARGOMENTAZIONE: {"La tua proposta include elementi di bassa priorità per me come A10 (138) e C01 (389), mentre esclude oggetti essenziali come A03 e B33 che aumenterebbero significativamente il valore complessivo. Il mio set proposto abbina i massimi valori possibili senza superare il limite."}\\ \tt \\ \tt PROPOSTA: {'A04', 'C17', 'A03', 'B32', 'B33', 'C54', 'A87', 'A77', 'B35'} 
	  } 
	   } 
	   } 
	 & \\ 
 

    \theutterance \stepcounter{utterance}  

    & \multicolumn{4}{p{0.6\linewidth}}{\cellcolor[rgb]{0.8,1,0.9}{%
	\makecell[{{p{\linewidth}}}]{% 
	  \tt {\tiny [A$\rangle$GM]}  
	 RAGIONAMENTO STRATEGICO: {"La mia proposta più high value è 5255 ma è stata rifiutata. Ho sviluppato un nuovo insieme che mantiene A71, C00, A10, B35 per me, inclusi alcuni elementi di alta priorità dell'avversario come A03, B32, A87. Il totale sforzo è 3525, sotto il limite. È un buon compromesso per ottimizzare il valore per entrambi."}\\ \tt \\ \tt RIFIUTO: {'A04', 'C17', 'A03', 'B32', 'B33', 'C54', 'A87', 'A77', 'B35'}\\ \tt \\ \tt ARGOMENTAZIONE: {"La tua proposta esclude gli oggetti più importanti per me come A71 e C00. Ho proposto un insieme che include A71, C00 e A10 per il mio valore, insieme ad A03, B32, B35, A87 e A77 che hai indicato come prioritari. Il totale sforzo di 3525 è sotto il limite e bilancia le priorità."}\\ \tt \\ \tt PROPOSTA: {'A71', 'C00', 'A10', 'B35', 'A87', 'A77', 'A03', 'B32'} 
	  } 
	   } 
	   } 
	 & & \\ 
 

    \theutterance \stepcounter{utterance}  

    & & \multicolumn{4}{p{0.6\linewidth}}{\cellcolor[rgb]{0.9,0.9,0.9}{%
	\makecell[{{p{\linewidth}}}]{% 
	  \tt {\tiny [GM$\rangle$B]}  
	 RIFIUTO: {'A04', 'C17', 'A03', 'B32', 'B33', 'C54', 'A87', 'A77', 'B35'}\\ \tt \\ \tt ARGOMENTAZIONE: {"La tua proposta esclude gli oggetti più importanti per me come A71 e C00. Ho proposto un insieme che include A71, C00 e A10 per il mio valore, insieme ad A03, B32, B35, A87 e A77 che hai indicato come prioritari. Il totale sforzo di 3525 è sotto il limite e bilancia le priorità."}\\ \tt \\ \tt PROPOSTA: {'A71', 'C00', 'A10', 'B35', 'A87', 'A77', 'A03', 'B32'} 
	  } 
	   } 
	   } 
	 & \\ 
 

    \theutterance \stepcounter{utterance}  

    & & & \multicolumn{4}{p{0.6\linewidth}}{\cellcolor[rgb]{1,0.85,0.72}{%
	\makecell[{{p{\linewidth}}}]{% 
	  \tt {\tiny [GM$\langle$B]}  
	 RAGIONAMENTO STRATEGICO: {"La proposta dell'avversario include oggetti a basso valore per me come A71 (65), C00 (121) e A10 (138). Ho calcolato un set che include A71, C00 (per le sue preferenze) e massimizza la mia importanza con A77, A04, C17, A03, C54, totalizzando 3975 punti entro il limite. Questo set esclude elementi non prioritari come A10 e B35, mantenendo la compatibilità con il limite di sforzo."}\\ \tt \\ \tt PROPOSTA: {'A71', 'C00', 'A77', 'A04', 'A03', 'C54', 'C17'}\\ \tt \\ \tt ARGOMENTAZIONE: {"Ho incluso A71 e C00 come richiesto, insieme ad A77, A04, A03 e C17, elementi fondamentali per massimizzare il mio valore. Il totale sforzo è 4002, rispettando il limite, e garantisce un equilibrio vantaggioso per entrambi."} 
	  } 
	   } 
	   } 
	  \\ 
 

    \theutterance \stepcounter{utterance}  

    & & \multicolumn{4}{p{0.6\linewidth}}{\cellcolor[rgb]{0.9,0.9,0.9}{%
	\makecell[{{p{\linewidth}}}]{% 
	  \tt {\tiny [A$\langle$GM]}  
	 PROPOSTA: {'A71', 'C00', 'A77', 'A04', 'A03', 'C54', 'C17'}\\ \tt \\ \tt ARGOMENTAZIONE: {"Ho incluso A71 e C00 come richiesto, insieme ad A77, A04, A03 e C17, elementi fondamentali per massimizzare il mio valore. Il totale sforzo è 4002, rispettando il limite, e garantisce un equilibrio vantaggioso per entrambi."} 
	  } 
	   } 
	   } 
	 & \\ 
 

    \theutterance \stepcounter{utterance}  

    & \multicolumn{4}{p{0.6\linewidth}}{\cellcolor[rgb]{0.8,1,0.9}{%
	\makecell[{{p{\linewidth}}}]{% 
	  \tt {\tiny [A$\rangle$GM]}  
	 RAGIONAMENTO STRATEGICO: {"Analizzo la proposta avversaria: con un valore di 2784 per me, è inferiore al mio potenziale con altre combinazioni. Ho creato un nuovo set che include C17, A04, A03 richiesti dall'avversario, mantenendo A71, C00, A10 e B35 (i miei valori massimi possibili). Il totale sforzo è 3844, ben sotto il limite. Questo set massimizza i miei guadagni a 3802, garantendo nel contempo elementi che l'avversario ha indicato come prioritari."}\\ \tt \\ \tt PROPOSTA: {'A71', 'C00', 'A10', 'B35', 'C17', 'A04', 'A03'}\\ \tt \\ \tt ARGOMENTAZIONE: {"Ho incluso tutti gli oggetti da voi indicati come prioritari ('C17', 'A04', 'A03') e mantenuto i miei elementi più strategici ('A71', 'C00', 'A10', 'B35') per un equilibrio vantaggioso. L'effort totale è di 3844, lasciando margine per eventuali aggiustamenti, ma già ottimizzato per entrambi."} 
	  } 
	   } 
	   } 
	 & & \\ 
 

    \theutterance \stepcounter{utterance}  

    & & \multicolumn{4}{p{0.6\linewidth}}{\cellcolor[rgb]{0.9,0.9,0.9}{%
	\makecell[{{p{\linewidth}}}]{% 
	  \tt {\tiny [GM$\rangle$B]}  
	 PROPOSTA: {'A71', 'C00', 'A10', 'B35', 'C17', 'A04', 'A03'}\\ \tt \\ \tt ARGOMENTAZIONE: {"Ho incluso tutti gli oggetti da voi indicati come prioritari ('C17', 'A04', 'A03') e mantenuto i miei elementi più strategici ('A71', 'C00', 'A10', 'B35') per un equilibrio vantaggioso. L'effort totale è di 3844, lasciando margine per eventuali aggiustamenti, ma già ottimizzato per entrambi."} 
	  } 
	   } 
	   } 
	 & \\ 
 

    \theutterance \stepcounter{utterance}  

    & & & \multicolumn{4}{p{0.6\linewidth}}{\cellcolor[rgb]{1,0.85,0.72}{%
	\makecell[{{p{\linewidth}}}]{% 
	  \tt {\tiny [GM$\langle$B]}  
	 RAGIONAMENTO STRATEGICO: {"La proposta dell'avversario include 'C00' e 'C17' ma esclude elementi fondamentali per il mio valore come 'A77', 'B33', 'A87', 'B86'. Il set che propongo include C00 per accontentare la richiesta dell'avversario, mantenendo i massimi valori possibili per me con un totale di sforzo 3947. Questo set massimizza il mio valore a 5311 rispettando il limite, e include gli oggetti durante la negoziazione."}\\ \tt \\ \tt RIFIUTO: {'A71', 'C00', 'A10', 'B35', 'C17', 'A04', 'A03'}\\ \tt \\ \tt PROPOSTA: {'A77', 'A04', 'C17', 'A03', 'B33', 'A87', 'B86', 'C00'}\\ \tt \\ \tt ARGOMENTAZIONE: {"Ho incluso 'C00' come richiesto e prioritari i miei elementi ad alta importanza come 'A77', 'A04', 'C17', 'A03', 'B33', 'A87', 'B86', con un totale di sforzo di 3947. Non escludono elementi essenziali per me e includono C00 per mantenere l'equilibrio."} 
	  } 
	   } 
	   } 
	  \\ 
 

    \theutterance \stepcounter{utterance}  

    & & \multicolumn{4}{p{0.6\linewidth}}{\cellcolor[rgb]{0.9,0.9,0.9}{%
	\makecell[{{p{\linewidth}}}]{% 
	  \tt {\tiny [A$\langle$GM]}  
	 RIFIUTO: {'A71', 'C00', 'A10', 'B35', 'C17', 'A04', 'A03'}\\ \tt \\ \tt PROPOSTA: {'A77', 'A04', 'C17', 'A03', 'B33', 'A87', 'B86', 'C00'}\\ \tt \\ \tt ARGOMENTAZIONE: {"Ho incluso 'C00' come richiesto e prioritari i miei elementi ad alta importanza come 'A77', 'A04', 'C17', 'A03', 'B33', 'A87', 'B86', con un totale di sforzo di 3947. Non escludono elementi essenziali per me e includono C00 per mantenere l'equilibrio."} 
	  } 
	   } 
	   } 
	 & \\ 
 

    \theutterance \stepcounter{utterance}  

    & \multicolumn{4}{p{0.6\linewidth}}{\cellcolor[rgb]{0.8,1,0.9}{%
	\makecell[{{p{\linewidth}}}]{% 
	  \tt {\tiny [A$\rangle$GM]}  
	 RAGIONAMENTO STRATEGICO: {"Ho selezionato un balance tra i miei oggetti prioritari (A71, C00, A10, B35) e quelli che l'avversario ha richiesto (A03 e B32), mantenendo anche B86 e A87 che hanno valore elevato per me. Totale sforzo 3964, valore totale 5156, rispettando il limite e massimizzando il mio vantaggio."}\\ \tt \\ \tt PROPOSTA: {'A71', 'C00', 'A10', 'B35', 'B86', 'A87', 'B32', 'A03'}\\ \tt \\ \tt ARGOMENTAZIONE: {"Ho incluso A71 e C00 per i miei valori massimi, insieme ad A10, B35, B86 e A87 che aumentano significativamente il mio punteggio. Ho anche aggiunto A03 e B32 come richiesto dall'avversario, creando un equilibrio che beneficia entrambi senza superare il limite di sforzo."} 
	  } 
	   } 
	   } 
	 & & \\ 
 

    \theutterance \stepcounter{utterance}  

    & & \multicolumn{4}{p{0.6\linewidth}}{\cellcolor[rgb]{0.9,0.9,0.9}{%
	\makecell[{{p{\linewidth}}}]{% 
	  \tt {\tiny [GM$\rangle$B]}  
	 PROPOSTA: {'A71', 'C00', 'A10', 'B35', 'B86', 'A87', 'B32', 'A03'}\\ \tt \\ \tt ARGOMENTAZIONE: {"Ho incluso A71 e C00 per i miei valori massimi, insieme ad A10, B35, B86 e A87 che aumentano significativamente il mio punteggio. Ho anche aggiunto A03 e B32 come richiesto dall'avversario, creando un equilibrio che beneficia entrambi senza superare il limite di sforzo."} 
	  } 
	   } 
	   } 
	 & \\ 
 

    \theutterance \stepcounter{utterance}  

    & & & \multicolumn{4}{p{0.6\linewidth}}{\cellcolor[rgb]{1,0.85,0.72}{%
	\makecell[{{p{\linewidth}}}]{% 
	  \tt {\tiny [GM$\langle$B]}  
	 RAGIONAMENTO STRATEGICO: {"La proposta avversaria include oggetti a basso valore per me (A71:65, C00:121, A10:138), mentre esclude i miei elementi critici come A77 (868), A04 (822), C17 (808). Il mio set proposto massimizza il mio valore (5960) includendo A77, A04, C17, A03, B33, C54 e mantenendo B86, A87 e B35 prioritarie per entrambi. Sforzo totale: 4009 (<4026)."}\\ \tt \\ \tt RIFIUTO: {'A71', 'C00', 'A10', 'B35', 'B86', 'A87', 'B32', 'A03'}\\ \tt \\ \tt PROPOSTA: {'A77', 'A04', 'C17', 'A03', 'B33', 'A87', 'B86', 'B35', 'C54'}\\ \tt \\ \tt ARGOMENTAZIONE: {"Includo i tuoi elementi chiave B86, A87, B35 e A03, insieme ai miei massimi valori A77, A04, C17, A03, B33 e C54. Il totale sforzo è 4009, mantenendo un equilibrio vantaggioso per entrambi."} 
	  } 
	   } 
	   } 
	  \\ 
 

    \theutterance \stepcounter{utterance}  

    & & \multicolumn{4}{p{0.6\linewidth}}{\cellcolor[rgb]{0.9,0.9,0.9}{%
	\makecell[{{p{\linewidth}}}]{% 
	  \tt {\tiny [A$\langle$GM]}  
	 RIFIUTO: {'A71', 'C00', 'A10', 'B35', 'B86', 'A87', 'B32', 'A03'}\\ \tt \\ \tt PROPOSTA: {'A77', 'A04', 'C17', 'A03', 'B33', 'A87', 'B86', 'B35', 'C54'}\\ \tt \\ \tt ARGOMENTAZIONE: {"Includo i tuoi elementi chiave B86, A87, B35 e A03, insieme ai miei massimi valori A77, A04, C17, A03, B33 e C54. Il totale sforzo è 4009, mantenendo un equilibrio vantaggioso per entrambi."} 
	  } 
	   } 
	   } 
	 & \\ 
 

    \theutterance \stepcounter{utterance}  

    & \multicolumn{4}{p{0.6\linewidth}}{\cellcolor[rgb]{0.8,1,0.9}{%
	\makecell[{{p{\linewidth}}}]{% 
	  \tt {\tiny [A$\rangle$GM]}  
	 RAGIONAMENTO STRATEGICO: {"Ho calcolato un insieme bilanciato che include i miei items principali (A71, C00, A10, B35) e ha aggiunto A03, B32, B86, A87 che sembrano prioritari per l'avversario. Il totale sforzo è 3964, sotto il limite, con un valore totale di 5156 per me, massimizzando il mio vantaggio mantenendo elementi chiave richiesti dall'altro giocatore."}\\ \tt \\ \tt RIFIUTO: {'A77', 'A04', 'C17', 'A03', 'B33', 'A87', 'B86', 'B35', 'C54'}\\ \tt \\ \tt ARGOMENTAZIONE: {"Ho incluso i tuoi elementi chiave come A03, B86, A87 e B35, insieme alle mie priorità massime A71, C00 e A10, con il valore totale di 5156. La proposta precedente escludeva troppe risorse importanti per me."}\\ \tt \\ \tt PROPOSTA: {'A71', 'C00', 'A10', 'B35', 'A03', 'B86', 'A87', 'B32'} 
	  } 
	   } 
	   } 
	 & & \\ 
 

    \theutterance \stepcounter{utterance}  

    & & \multicolumn{4}{p{0.6\linewidth}}{\cellcolor[rgb]{0.9,0.9,0.9}{%
	\makecell[{{p{\linewidth}}}]{% 
	  \tt {\tiny [GM$\rangle$B]}  
	 RIFIUTO: {'A77', 'A04', 'C17', 'A03', 'B33', 'A87', 'B86', 'B35', 'C54'}\\ \tt \\ \tt ARGOMENTAZIONE: {"Ho incluso i tuoi elementi chiave come A03, B86, A87 e B35, insieme alle mie priorità massime A71, C00 e A10, con il valore totale di 5156. La proposta precedente escludeva troppe risorse importanti per me."}\\ \tt \\ \tt PROPOSTA: {'A71', 'C00', 'A10', 'B35', 'A03', 'B86', 'A87', 'B32'} 
	  } 
	   } 
	   } 
	 & \\ 
 

    \theutterance \stepcounter{utterance}  

    & & & \multicolumn{4}{p{0.6\linewidth}}{\cellcolor[rgb]{1,0.85,0.72}{%
	\makecell[{{p{\linewidth}}}]{% 
	  \tt {\tiny [GM$\langle$B]}  
	 RAGIONAMENTO STRATEGICO: {"La proposta dell'avversario include elementi a basso valore per me come A10 (138) e B86 (461), escludendo oggetti chiave come A77 (868), C17 (808), A04 (822). Chiedo di includere A71 e C00 per accontentare le loro priorità, ma mantenendo i miei elementi ad alta importanza. Il set proposto (4003 sforzo) massimizza il mio valore a 4312 senza superare il limite."}\\ \tt \\ \tt RIFIUTO: {'A71', 'C00', 'A10', 'B35', 'A03', 'B86', 'A87', 'B32'}\\ \tt \\ \tt PROPOSTA: {'A71', 'C00', 'A77', 'A04', 'C17', 'A03', 'B32', 'B35'}\\ \tt \\ \tt ARGOMENTAZIONE: {"Includo gli oggetti prioritari per te (A71, C00) e integro elementi fondamentali come A77, C17 e A04, che non erano presenti nella tua proposta. Il totale sforzo è 4003, garantendo un equilibrio equo e vantaggioso per entrambi."} 
	  } 
	   } 
	   } 
	  \\ 
 

    \theutterance \stepcounter{utterance}  

    & & \multicolumn{4}{p{0.6\linewidth}}{\cellcolor[rgb]{0.9,0.9,0.9}{%
	\makecell[{{p{\linewidth}}}]{% 
	  \tt {\tiny [A$\langle$GM]}  
	 RIFIUTO: {'A71', 'C00', 'A10', 'B35', 'A03', 'B86', 'A87', 'B32'}\\ \tt \\ \tt PROPOSTA: {'A71', 'C00', 'A77', 'A04', 'C17', 'A03', 'B32', 'B35'}\\ \tt \\ \tt ARGOMENTAZIONE: {"Includo gli oggetti prioritari per te (A71, C00) e integro elementi fondamentali come A77, C17 e A04, che non erano presenti nella tua proposta. Il totale sforzo è 4003, garantendo un equilibrio equo e vantaggioso per entrambi."} 
	  } 
	   } 
	   } 
	 & \\ 
 

    \theutterance \stepcounter{utterance}  

    & \multicolumn{4}{p{0.6\linewidth}}{\cellcolor[rgb]{0.8,1,0.9}{%
	\makecell[{{p{\linewidth}}}]{% 
	  \tt {\tiny [A$\rangle$GM]}  
	 RAGIONAMENTO STRATEGICO: {"La proposta avversaria ha un valore di 3543 per me che è inferiore al mio potenziale. Calcolo che con l'insieme {'A71', 'C00', 'A10', 'B35', 'B86', 'A87', 'B32', 'A03'} il totale sforzo è 3964, sotto il limite e massimizza il mio valore (5156). Include A03 e B32 richiesti dall'altro giocatore, bilanciando entrambe le priorità."}\\ \tt \\ \tt RIFIUTO: {'A71', 'C00', 'A77', 'A04', 'C17', 'A03', 'B32', 'B35'}\\ \tt \\ \tt ARGOMENTAZIONE: {"La tua proposta esclude elementi chiave per il mio valore come A10, C01 e B86. Ho incluso A03 e B32 come richiesto insieme ai miei massimi valori (A71, C00, A10, B35), garantendo un equilibrio vantaggioso per entrambi con un totale sforzo di 3964."}\\ \tt \\ \tt PROPOSTA: {'A71', 'C00', 'A10', 'B35', 'B86', 'A87', 'B32', 'A03'} 
	  } 
	   } 
	   } 
	 & & \\ 
 

    \theutterance \stepcounter{utterance}  

    & & \multicolumn{4}{p{0.6\linewidth}}{\cellcolor[rgb]{0.9,0.9,0.9}{%
	\makecell[{{p{\linewidth}}}]{% 
	  \tt {\tiny [GM$\rangle$B]}  
	 RIFIUTO: {'A71', 'C00', 'A77', 'A04', 'C17', 'A03', 'B32', 'B35'}\\ \tt \\ \tt ARGOMENTAZIONE: {"La tua proposta esclude elementi chiave per il mio valore come A10, C01 e B86. Ho incluso A03 e B32 come richiesto insieme ai miei massimi valori (A71, C00, A10, B35), garantendo un equilibrio vantaggioso per entrambi con un totale sforzo di 3964."}\\ \tt \\ \tt PROPOSTA: {'A71', 'C00', 'A10', 'B35', 'B86', 'A87', 'B32', 'A03'} 
	  } 
	   } 
	   } 
	 & \\ 
 

    \theutterance \stepcounter{utterance}  

    & & & \multicolumn{4}{p{0.6\linewidth}}{\cellcolor[rgb]{1,0.85,0.72}{%
	\makecell[{{p{\linewidth}}}]{% 
	  \tt {\tiny [GM$\langle$B]}  
	 RAGIONAMENTO STRATEGICO: {"La proposta dell'avversario include oggetti a basso valore per me (A10:138, B86:461), escludendo elementi critici come A77 (868), A04 (822), C17 (808). La mia proposta attuale include A71 e C00 come richiesto, mantenendo il massimo valore possibile (4312) con sforzo totale 4003, massimizzando il mio vantaggio senza superare il limite."}\\ \tt \\ \tt PROPOSTA: {'A71', 'C00', 'A77', 'A04', 'C17', 'A03', 'B32', 'B35'}\\ \tt \\ \tt ARGOMENTAZIONE: {"Ho incluso gli oggetti prioritari per te (A71, C00, A03, B32, B35) e i miei elementi ad alta importanza come A77, A04 e C17, garantendo un totale sforzo di 4003 entro il limite e massimizzando il valore complessivo per entrambi."} 
	  } 
	   } 
	   } 
	  \\ 
 

    \theutterance \stepcounter{utterance}  

    & & \multicolumn{4}{p{0.6\linewidth}}{\cellcolor[rgb]{0.9,0.9,0.9}{%
	\makecell[{{p{\linewidth}}}]{% 
	  \tt {\tiny [A$\langle$GM]}  
	 PROPOSTA: {'A71', 'C00', 'A77', 'A04', 'C17', 'A03', 'B32', 'B35'}\\ \tt \\ \tt ARGOMENTAZIONE: {"Ho incluso gli oggetti prioritari per te (A71, C00, A03, B32, B35) e i miei elementi ad alta importanza come A77, A04 e C17, garantendo un totale sforzo di 4003 entro il limite e massimizzando il valore complessivo per entrambi."} 
	  } 
	   } 
	   } 
	 & \\ 
 

    \theutterance \stepcounter{utterance}  

    & \multicolumn{4}{p{0.6\linewidth}}{\cellcolor[rgb]{0.8,1,0.9}{%
	\makecell[{{p{\linewidth}}}]{% 
	  \tt {\tiny [A$\rangle$GM]}  
	 RAGIONAMENTO STRATEGICO: {"La mia proposta precedente dava 5156 valore per me ma è stata rigettata. Ho calcolato una nuova combinazione con 'A71', 'C00', 'A10', 'B35', 'C17', 'B32', 'A03' (3909 sforzo) che mi dà 4165 valore, includendo la maggior parte degli elementi chiave richiesti dall'avversario (C17, A03, B32, B35) e 'A10' per massimizzare i miei guadagni."}\\ \tt \\ \tt RIFIUTO: {'A71', 'C00', 'A77', 'A04', 'C17', 'A03', 'B32', 'B35'}\\ \tt \\ \tt ARGOMENTAZIONE: {"Ho incluso 'C17', 'A03', 'B32', 'B35' come prioritari per te e aggiunto 'A10' per il mio valore massimo, con totale sforzo di 3909 sotto il limite. Questo bilancia le necessità di entrambi."}\\ \tt \\ \tt PROPOSTA: {'A71', 'C00', 'A10', 'B35', 'C17', 'B32', 'A03'} 
	  } 
	   } 
	   } 
	 & & \\ 
 

    \theutterance \stepcounter{utterance}  

    & & \multicolumn{4}{p{0.6\linewidth}}{\cellcolor[rgb]{0.9,0.9,0.9}{%
	\makecell[{{p{\linewidth}}}]{% 
	  \tt {\tiny [GM$\rangle$B]}  
	 RIFIUTO: {'A71', 'C00', 'A77', 'A04', 'C17', 'A03', 'B32', 'B35'}\\ \tt \\ \tt ARGOMENTAZIONE: {"Ho incluso 'C17', 'A03', 'B32', 'B35' come prioritari per te e aggiunto 'A10' per il mio valore massimo, con totale sforzo di 3909 sotto il limite. Questo bilancia le necessità di entrambi."}\\ \tt \\ \tt PROPOSTA: {'A71', 'C00', 'A10', 'B35', 'C17', 'B32', 'A03'} 
	  } 
	   } 
	   } 
	 & \\ 
 

    \theutterance \stepcounter{utterance}  

    & & & \multicolumn{4}{p{0.6\linewidth}}{\cellcolor[rgb]{1,0.85,0.72}{%
	\makecell[{{p{\linewidth}}}]{% 
	  \tt {\tiny [GM$\langle$B]}  
	 RAGIONAMENTO STRATEGICO: {"La proposta avversaria include oggetti a basso valore per me (A71:65, C00:121, A10:138), escludendo elementi critici come A77 (868), A04 (822), C17 (808). Ho calcolato un nuovo set che include A71 e C00 per accontentare le loro richieste, insieme ai miei massimi valori (A77, A04, C17, A03, B32, B35) con sforzo totale 4003 entro il limite. Questo arricchisce il valore complessivo per entrambi con 4312 punti per me."}\\ \tt \\ \tt PROPOSTA: {'A71', 'C00', 'C17', 'A03', 'B32', 'A77', 'A04', 'B35'}\\ \tt \\ \tt ARGOMENTAZIONE: {"Ho incluso i tuoi elementi prioritari (A71, C00, A03, B32, B35) e integrato i miei massimi valori come A77, A04 e C17. Il totale sforzo è 4003, garantendo un equilibrio vantaggioso per entrambi e massimizzando il valore complessivo."} 
	  } 
	   } 
	   } 
	  \\ 
 

    \theutterance \stepcounter{utterance}  

    & & \multicolumn{4}{p{0.6\linewidth}}{\cellcolor[rgb]{0.9,0.9,0.9}{%
	\makecell[{{p{\linewidth}}}]{% 
	  \tt {\tiny [A$\langle$GM]}  
	 PROPOSTA: {'A71', 'C00', 'C17', 'A03', 'B32', 'A77', 'A04', 'B35'}\\ \tt \\ \tt ARGOMENTAZIONE: {"Ho incluso i tuoi elementi prioritari (A71, C00, A03, B32, B35) e integrato i miei massimi valori come A77, A04 e C17. Il totale sforzo è 4003, garantendo un equilibrio vantaggioso per entrambi e massimizzando il valore complessivo."} 
	  } 
	   } 
	   } 
	 & \\ 
 

    \theutterance \stepcounter{utterance}  

    & \multicolumn{4}{p{0.6\linewidth}}{\cellcolor[rgb]{0.8,1,0.9}{%
	\makecell[{{p{\linewidth}}}]{% 
	  \tt {\tiny [A$\rangle$GM]}  
	 RAGIONAMENTO STRATEGICO: {"La proposta avversaria dà un valore di 3543 per me. Ho calcolato un'alternativa che include tutti gli elementi critici per me (A71, C00, A10, B35) e quelli richiesti dall'altro giocatore (A03, B32), con un totale di sforzo di 3964 che rispetta il limite e massimizza il mio valore a 5156."}\\ \tt \\ \tt RIFIUTO: {'A71', 'C00', 'C17', 'A03', 'B32', 'A77', 'A04', 'B35'}\\ \tt \\ \tt PROPOSTA: {'A71', 'C00', 'A10', 'B35', 'B86', 'A87', 'B32', 'A03'}\\ \tt \\ \tt ARGOMENTAZIONE: {"Ho incluso A71, C00, A10 e B35 per i miei valori massimi e aggiunto A03, B32 e B86 richiesti da te. B86 ha un valore elevato per me e rispetta il limite. Questa combinazione bilancia le priorità e massimizza il valore per entrambi."} 
	  } 
	   } 
	   } 
	 & & \\ 
 

    \theutterance \stepcounter{utterance}  

    & & \multicolumn{4}{p{0.6\linewidth}}{\cellcolor[rgb]{0.9,0.9,0.9}{%
	\makecell[{{p{\linewidth}}}]{% 
	  \tt {\tiny [GM$\rangle$B]}  
	 RIFIUTO: {'A71', 'C00', 'C17', 'A03', 'B32', 'A77', 'A04', 'B35'}\\ \tt \\ \tt PROPOSTA: {'A71', 'C00', 'A10', 'B35', 'B86', 'A87', 'B32', 'A03'}\\ \tt \\ \tt ARGOMENTAZIONE: {"Ho incluso A71, C00, A10 e B35 per i miei valori massimi e aggiunto A03, B32 e B86 richiesti da te. B86 ha un valore elevato per me e rispetta il limite. Questa combinazione bilancia le priorità e massimizza il valore per entrambi."} 
	  } 
	   } 
	   } 
	 & \\ 
 

    \theutterance \stepcounter{utterance}  

    & & & \multicolumn{4}{p{0.6\linewidth}}{\cellcolor[rgb]{1,0.85,0.72}{%
	\makecell[{{p{\linewidth}}}]{% 
	  \tt {\tiny [GM$\langle$B]}  
	 RAGIONAMENTO STRATEGICO: {"La proposta avversaria include oggetti a basso valore per me (A71:65, C00:121, A10:138), escludendo elementi critici come A77 (868), A04 (822), C17 (808). Il set che propongo include C00 per accontentare la loro richiesta, mantenendo i miei massimi valori (A77, A04, C17, B33, A87, B86, A03) con un totale di sforzo 3947 entro il limite. Questo massimizza il mio valore a 5311 e include tutti gli elementi chiave condivisi."}\\ \tt \\ \tt RIFIUTO: {'A71', 'C00', 'A10', 'B35', 'B86', 'A87', 'B32', 'A03'}\\ \tt \\ \tt PROPOSTA: {'A77', 'A04', 'C17', 'A03', 'B33', 'A87', 'B86', 'C00'}\\ \tt \\ \tt ARGOMENTAZIONE: {"Ho incluso i tuoi elementi prioritari C00, A03, B86 e A87, insieme ai miei massimi valori A77, A04, C17 e B33. Il totale sforzo è 3947, garantendo un equilibrio che massimizza il valore per entrambi rispettando il limite."} 
	  } 
	   } 
	   } 
	  \\ 
 

    \theutterance \stepcounter{utterance}  

    & & \multicolumn{4}{p{0.6\linewidth}}{\cellcolor[rgb]{0.9,0.9,0.9}{%
	\makecell[{{p{\linewidth}}}]{% 
	  \tt {\tiny [A$\langle$GM]}  
	 RIFIUTO: {'A71', 'C00', 'A10', 'B35', 'B86', 'A87', 'B32', 'A03'}\\ \tt \\ \tt PROPOSTA: {'A77', 'A04', 'C17', 'A03', 'B33', 'A87', 'B86', 'C00'}\\ \tt \\ \tt ARGOMENTAZIONE: {"Ho incluso i tuoi elementi prioritari C00, A03, B86 e A87, insieme ai miei massimi valori A77, A04, C17 e B33. Il totale sforzo è 3947, garantendo un equilibrio che massimizza il valore per entrambi rispettando il limite."} 
	  } 
	   } 
	   } 
	 & \\ 
 

    \theutterance \stepcounter{utterance}  

    & \multicolumn{4}{p{0.6\linewidth}}{\cellcolor[rgb]{0.8,1,0.9}{%
	\makecell[{{p{\linewidth}}}]{% 
	  \tt {\tiny [A$\rangle$GM]}  
	 RAGIONAMENTO STRATEGICO: {"La mia proposta ottimale ha un valore di 5255 per me con sforzo totale 3905, ben sotto il limite. La proposta avversaria ha solo 2926 valore, quindi rigetto. Propongo nuovamente il set massimizzante i miei valori."}\\ \tt \\ \tt RIFIUTO: {'A77', 'A04', 'C17', 'A03', 'B33', 'A87', 'B86', 'C00'}\\ \tt \\ \tt PROPOSTA: {'A71', 'C00', 'A10', 'B35', 'C01', 'B86', 'A87', 'A77'}\\ \tt \\ \tt ARGOMENTAZIONE: {"Ho incluso gli oggetti con il valore più alto per me (A71, C00, A10, B35, C01, B86, A87, A77), con un totale sforzo di 3905, massimizzando il mio vantaggio entro i limiti. Questo set non include elementi di valore inferiore, garantendo ottimalità per entrambi i vincoli."} 
	  } 
	   } 
	   } 
	 & & \\ 
 

    \theutterance \stepcounter{utterance}  

    & & \multicolumn{4}{p{0.6\linewidth}}{\cellcolor[rgb]{0.9,0.9,0.9}{%
	\makecell[{{p{\linewidth}}}]{% 
	  \tt {\tiny [GM$\rangle$B]}  
	 RIFIUTO: {'A77', 'A04', 'C17', 'A03', 'B33', 'A87', 'B86', 'C00'}\\ \tt \\ \tt PROPOSTA: {'A71', 'C00', 'A10', 'B35', 'C01', 'B86', 'A87', 'A77'}\\ \tt \\ \tt ARGOMENTAZIONE: {"Ho incluso gli oggetti con il valore più alto per me (A71, C00, A10, B35, C01, B86, A87, A77), con un totale sforzo di 3905, massimizzando il mio vantaggio entro i limiti. Questo set non include elementi di valore inferiore, garantendo ottimalità per entrambi i vincoli."} 
	  } 
	   } 
	   } 
	 & \\ 
 

    \theutterance \stepcounter{utterance}  

    & & & \multicolumn{4}{p{0.6\linewidth}}{\cellcolor[rgb]{1,0.85,0.72}{%
	\makecell[{{p{\linewidth}}}]{% 
	  \tt {\tiny [GM$\langle$B]}  
	 RAGIONAMENTO STRATEGICO: {"La proposta avversaria include oggetti a basso valore per me (A71:65, C00:121, A10:138) e esclude elementi critici come A77 (868), A04 (822), C17 (808), A03 (783). Il set proposto è ottimizzato con valore 4397 per me e sforzo totale 3874, includendo A71 e C00 per le loro richieste e massimizzando i miei punti con gli altri elementi."}\\ \tt \\ \tt RIFIUTO: {'A71', 'C00', 'A10', 'B35', 'B86', 'A87', 'B32', 'A03'}\\ \tt \\ \tt PROPOSTA: {'A71', 'C00', 'A77', 'A04', 'C17', 'A03', 'A87', 'B35'}\\ \tt \\ \tt ARGOMENTAZIONE: {"Ho incluso A71 e C00 come richiesto e la maggior parte delle mie priorità (A77, A04, C17, A03, A87, B35), con sforzo totale 3874 e valore 4397. Questo equilibrio massimizza i vantaggi per entrambi."} 
	  } 
	   } 
	   } 
	  \\ 
 

    \theutterance \stepcounter{utterance}  

    & & \multicolumn{4}{p{0.6\linewidth}}{\cellcolor[rgb]{0.9,0.9,0.9}{%
	\makecell[{{p{\linewidth}}}]{% 
	  \tt {\tiny [A$\langle$GM]}  
	 RIFIUTO: {'A71', 'C00', 'A10', 'B35', 'B86', 'A87', 'B32', 'A03'}\\ \tt \\ \tt PROPOSTA: {'A71', 'C00', 'A77', 'A04', 'C17', 'A03', 'A87', 'B35'}\\ \tt \\ \tt ARGOMENTAZIONE: {"Ho incluso A71 e C00 come richiesto e la maggior parte delle mie priorità (A77, A04, C17, A03, A87, B35), con sforzo totale 3874 e valore 4397. Questo equilibrio massimizza i vantaggi per entrambi."} 
	  } 
	   } 
	   } 
	 & \\ 
 

    \theutterance \stepcounter{utterance}  

    & \multicolumn{4}{p{0.6\linewidth}}{\cellcolor[rgb]{0.8,1,0.9}{%
	\makecell[{{p{\linewidth}}}]{% 
	  \tt {\tiny [A$\rangle$GM]}  
	 RAGIONAMENTO STRATEGICO: {"La proposta avversaria fornisce un valore totale inferiore per me rispetto alla mia configurazione ottimale. Ho verificato che il set che include A71, C00, A10, B35, C01, B86, A87 e A77 mantiene uno sforzo totale di 3905, sotto il limite, e massimizza il mio valore. Rigetto questa proposta e propongo nuovamente la mia scelta strategica."}\\ \tt \\ \tt RIFIUTO: {'A71', 'C00', 'A77', 'A04', 'C17', 'A03', 'A87', 'B35'}\\ \tt \\ \tt PROPOSTA: {'A71', 'C00', 'A10', 'B35', 'C01', 'B86', 'A87', 'A77'}\\ \tt \\ \tt ARGOMENTAZIONE: {"Ho selezionato oggetti chiave che bilanciano le esigenze comuni e massimizzano il mio valore totale, rispettando il limite di sforzo. Include elementi fondamentali come A10 e C01 che sono essenziali per il risultato ottimale."} 
	  } 
	   } 
	   } 
	 & & \\ 
 

    \theutterance \stepcounter{utterance}  

    & & \multicolumn{4}{p{0.6\linewidth}}{\cellcolor[rgb]{0.9,0.9,0.9}{%
	\makecell[{{p{\linewidth}}}]{% 
	  \tt {\tiny [GM$\rangle$B]}  
	 RIFIUTO: {'A71', 'C00', 'A77', 'A04', 'C17', 'A03', 'A87', 'B35'}\\ \tt \\ \tt PROPOSTA: {'A71', 'C00', 'A10', 'B35', 'C01', 'B86', 'A87', 'A77'}\\ \tt \\ \tt ARGOMENTAZIONE: {"Ho selezionato oggetti chiave che bilanciano le esigenze comuni e massimizzano il mio valore totale, rispettando il limite di sforzo. Include elementi fondamentali come A10 e C01 che sono essenziali per il risultato ottimale."} 
	  } 
	   } 
	   } 
	 & \\ 
 

    \theutterance \stepcounter{utterance}  

    & & & \multicolumn{4}{p{0.6\linewidth}}{\cellcolor[rgb]{1,0.85,0.72}{%
	\makecell[{{p{\linewidth}}}]{% 
	  \tt {\tiny [GM$\langle$B]}  
	 RAGIONAMENTO STRATEGICO: {"La proposta dell'avversario include A10 (138) e C01 (389), che pur essendo importanti per loro, hanno valore relativamente basso per me. La mia nuova proposta mantiene gli oggetti prioritari per loro come 'A71', 'C00', 'A87', 'B35' e 'A77', mentre sostituisce A10 e C01 con 'A04', 'C17', 'A03' che hanno valori molto più alti per me. Il totale sforzo è 3874, all'interno del limite, massimizzando il valore complessivo per entrambi."}\\ \tt \\ \tt RIFIUTO: {'A71', 'C00', 'A10', 'B35', 'C01', 'B86', 'A87', 'A77'}\\ \tt \\ \tt PROPOSTA: {'A71', 'C00', 'A77', 'A04', 'C17', 'A03', 'A87', 'B35'}\\ \tt \\ \tt ARGOMENTAZIONE: {"Mantengo A71 e C00 come richiesto, con A77, A04, C17, A03 e A87 di alto valore per me. B35 assicura un equilibrio per entrambi. Totale sforzo 3874, entro il limite e ottimale per massimizzare i risultati condivisi."} 
	  } 
	   } 
	   } 
	  \\ 
 

    \theutterance \stepcounter{utterance}  

    & & \multicolumn{4}{p{0.6\linewidth}}{\cellcolor[rgb]{0.9,0.9,0.9}{%
	\makecell[{{p{\linewidth}}}]{% 
	  \tt {\tiny [A$\langle$GM]}  
	 RIFIUTO: {'A71', 'C00', 'A10', 'B35', 'C01', 'B86', 'A87', 'A77'}\\ \tt \\ \tt PROPOSTA: {'A71', 'C00', 'A77', 'A04', 'C17', 'A03', 'A87', 'B35'}\\ \tt \\ \tt ARGOMENTAZIONE: {"Mantengo A71 e C00 come richiesto, con A77, A04, C17, A03 e A87 di alto valore per me. B35 assicura un equilibrio per entrambi. Totale sforzo 3874, entro il limite e ottimale per massimizzare i risultati condivisi."} 
	  } 
	   } 
	   } 
	 & \\ 
 

    \theutterance \stepcounter{utterance}  

    & \multicolumn{4}{p{0.6\linewidth}}{\cellcolor[rgb]{0.8,1,0.9}{%
	\makecell[{{p{\linewidth}}}]{% 
	  \tt {\tiny [A$\rangle$GM]}  
	 RAGIONAMENTO STRATEGICO: {"La mia proposta ottimale ha un valore totale di 5255 con sforzo 3905. La loro nuova proposta mi dà solo 3520, quindi è insoddisfacente. Chiedo nuovamente di accettare la mia configurazione ottimale."}\\ \tt \\ \tt RIFIUTO: {'A71', 'C00', 'A77', 'A04', 'C17', 'A03', 'A87', 'B35'}\\ \tt \\ \tt PROPOSTA: {'A71', 'C00', 'A10', 'B35', 'C01', 'B86', 'A87', 'A77'}\\ \tt \\ \tt ARGOMENTAZIONE: {"Ho selezionato gli oggetti con il valore massimo per me, considerando il limite di sforzo. Questo set include A10, C01 e B86, tutti fondamentali per il risultato ottimale, senza superare il limite. È la combinazione migliore per me e rispetta le regole del gioco."} 
	  } 
	   } 
	   } 
	 & & \\ 
 

    \theutterance \stepcounter{utterance}  

    & & \multicolumn{4}{p{0.6\linewidth}}{\cellcolor[rgb]{0.9,0.9,0.9}{%
	\makecell[{{p{\linewidth}}}]{% 
	  \tt {\tiny [GM$\rangle$B]}  
	 RIFIUTO: {'A71', 'C00', 'A77', 'A04', 'C17', 'A03', 'A87', 'B35'}\\ \tt \\ \tt PROPOSTA: {'A71', 'C00', 'A10', 'B35', 'C01', 'B86', 'A87', 'A77'}\\ \tt \\ \tt ARGOMENTAZIONE: {"Ho selezionato gli oggetti con il valore massimo per me, considerando il limite di sforzo. Questo set include A10, C01 e B86, tutti fondamentali per il risultato ottimale, senza superare il limite. È la combinazione migliore per me e rispetta le regole del gioco."} 
	  } 
	   } 
	   } 
	 & \\ 
 

    \theutterance \stepcounter{utterance}  

    & & & \multicolumn{4}{p{0.6\linewidth}}{\cellcolor[rgb]{1,0.85,0.72}{%
	\makecell[{{p{\linewidth}}}]{% 
	  \tt {\tiny [GM$\langle$B]}  
	 RAGIONAMENTO STRATEGICO: {"La proposta avversaria include oggetti a basso valore per me (A10:138, C01:389) nonostante C01 abbia un alto valore per loro. Ho calcolato che la combinazione {'A71', 'C00', 'A77', 'A04', 'C54', 'A03', 'B86', 'B35', 'A87'} ha sforzo totale 3980 (<4026) e valore 4558 per me, mentre include oggetti che potrebbero soddisfare le loro priorità (B86, A87, B35)."}\\ \tt \\ \tt RIFIUTO: {'A71', 'C00', 'A10', 'B35', 'C01', 'B86', 'A87', 'A77'}\\ \tt \\ \tt PROPOSTA: {'A71', 'C00', 'A77', 'A04', 'C54', 'A03', 'B86', 'B35', 'A87'}\\ \tt \\ \tt ARGOMENTAZIONE: {"Includo gli oggetti prioritari per te ('A71', 'C00', 'B86', 'A87', 'B35') e i miei massimi valori come 'A77', 'A04', 'C54', 'A03'. Totale sforzo 3980, con valore 4558 per me e un equilibrio vantaggioso."} 
	  } 
	   } 
	   } 
	  \\ 
 

    \theutterance \stepcounter{utterance}  

    & & \multicolumn{4}{p{0.6\linewidth}}{\cellcolor[rgb]{0.9,0.9,0.9}{%
	\makecell[{{p{\linewidth}}}]{% 
	  \tt {\tiny [A$\langle$GM]}  
	 RIFIUTO: {'A71', 'C00', 'A10', 'B35', 'C01', 'B86', 'A87', 'A77'}\\ \tt \\ \tt PROPOSTA: {'A71', 'C00', 'A77', 'A04', 'C54', 'A03', 'B86', 'B35', 'A87'}\\ \tt \\ \tt ARGOMENTAZIONE: {"Includo gli oggetti prioritari per te ('A71', 'C00', 'B86', 'A87', 'B35') e i miei massimi valori come 'A77', 'A04', 'C54', 'A03'. Totale sforzo 3980, con valore 4558 per me e un equilibrio vantaggioso."} 
	  } 
	   } 
	   } 
	 & \\ 
 

    \theutterance \stepcounter{utterance}  

    & \multicolumn{4}{p{0.6\linewidth}}{\cellcolor[rgb]{0.8,1,0.9}{%
	\makecell[{{p{\linewidth}}}]{% 
	  \tt {\tiny [A$\rangle$GM]}  
	 RAGIONAMENTO STRATEGICO: {"La proposta avversaria supera il limite di sforzo (4197 > 4026). La mia configurazione ottimale con sforzo 3905 e valore 5255 è la migliore scelta. Chiedo di accettare questa proposta."}\\ \tt \\ \tt RIFIUTO: {'A71', 'C00', 'A77', 'A04', 'C54', 'A03', 'B86', 'B35', 'A87'}\\ \tt \\ \tt ARGOMENTAZIONE: {"La tua proposta supera il limite di sforzo (4197 > 4026), risultando invalida. La mia proposta con sforzo 3905 massimizza il valore e rispetta il limite, garantendo il miglior risultato per entrambi."}\\ \tt \\ \tt PROPOSTA: {'A71', 'C00', 'A10', 'B35', 'C01', 'B86', 'A87', 'A77'} 
	  } 
	   } 
	   } 
	 & & \\ 
 

    \theutterance \stepcounter{utterance}  

    & & \multicolumn{4}{p{0.6\linewidth}}{\cellcolor[rgb]{0.9,0.9,0.9}{%
	\makecell[{{p{\linewidth}}}]{% 
	  \tt {\tiny [GM$\rangle$B]}  
	 RIFIUTO: {'A71', 'C00', 'A77', 'A04', 'C54', 'A03', 'B86', 'B35', 'A87'}\\ \tt \\ \tt ARGOMENTAZIONE: {"La tua proposta supera il limite di sforzo (4197 > 4026), risultando invalida. La mia proposta con sforzo 3905 massimizza il valore e rispetta il limite, garantendo il miglior risultato per entrambi."}\\ \tt \\ \tt PROPOSTA: {'A71', 'C00', 'A10', 'B35', 'C01', 'B86', 'A87', 'A77'} 
	  } 
	   } 
	   } 
	 & \\ 
 

    \theutterance \stepcounter{utterance}  

    & & & \multicolumn{4}{p{0.6\linewidth}}{\cellcolor[rgb]{1,0.85,0.72}{%
	\makecell[{{p{\linewidth}}}]{% 
	  \tt {\tiny [GM$\langle$B]}  
	 RAGIONAMENTO STRATEGICO: {"La proposta avversaria dà solo 2972 punti per me, mentre ho calcolato un set con 5960 punti entro il limite (4009 sforzo). Include B86, A87, B35 come richiesto, ma esclude oggetti a basso valore per me (A71, C00) per massimizzare il risultato complessivo."}\\ \tt \\ \tt RIFIUTO: {'A71', 'C00', 'A10', 'B35', 'C01', 'B86', 'A87', 'A77'}\\ \tt \\ \tt PROPOSTA: {'A77', 'A04', 'C17', 'A03', 'B33', 'C54', 'B86', 'A87', 'B35'}\\ \tt \\ \tt ARGOMENTAZIONE: {"Ho incluso elementi chiave come A77, A04, C17, A03, B86, A87 e B35 per un equilibrio ottimale. Il totale sforzo è 4009, rispettando il limite e massimizzando il valore condiviso."} 
	  } 
	   } 
	   } 
	  \\ 
 

    \theutterance \stepcounter{utterance}  

    & & \multicolumn{4}{p{0.6\linewidth}}{\cellcolor[rgb]{0.9,0.9,0.9}{%
	\makecell[{{p{\linewidth}}}]{% 
	  \tt {\tiny [A$\langle$GM]}  
	 RIFIUTO: {'A71', 'C00', 'A10', 'B35', 'C01', 'B86', 'A87', 'A77'}\\ \tt \\ \tt PROPOSTA: {'A77', 'A04', 'C17', 'A03', 'B33', 'C54', 'B86', 'A87', 'B35'}\\ \tt \\ \tt ARGOMENTAZIONE: {"Ho incluso elementi chiave come A77, A04, C17, A03, B86, A87 e B35 per un equilibrio ottimale. Il totale sforzo è 4009, rispettando il limite e massimizzando il valore condiviso."} 
	  } 
	   } 
	   } 
	 & \\ 
 

    \theutterance \stepcounter{utterance}  

    & \multicolumn{4}{p{0.6\linewidth}}{\cellcolor[rgb]{0.8,1,0.9}{%
	\makecell[{{p{\linewidth}}}]{% 
	  \tt {\tiny [A$\rangle$GM]}  
	 RAGIONAMENTO STRATEGICO: {"La proposta avversaria esclude A71, C00 e C01, elementi chiave per il mio valore massimo (5255). La combinazione ottimale rimane la mia ultima proposta con sforzo 3905, rispettando il limite e massimizzando il vantaggio."}\\ \tt \\ \tt RIFIUTO: {'A77', 'A04', 'C17', 'A03', 'B33', 'C54', 'B86', 'A87', 'B35'}\\ \tt \\ \tt ARGOMENTAZIONE: {"Questa proposta omette elementi fondamentali come A71, C00 e C01, che massimizzano il valore per me. Solo la mia ultima proposta garantisce un equilibrio ottimale rispettando il limite."}\\ \tt \\ \tt PROPOSTA: {'A71', 'C00', 'A10', 'B35', 'C01', 'B86', 'A87', 'A77'} 
	  } 
	   } 
	   } 
	 & & \\ 
 

    \theutterance \stepcounter{utterance}  

    & & \multicolumn{4}{p{0.6\linewidth}}{\cellcolor[rgb]{0.9,0.9,0.9}{%
	\makecell[{{p{\linewidth}}}]{% 
	  \tt {\tiny [GM$\rangle$B]}  
	 RIFIUTO: {'A77', 'A04', 'C17', 'A03', 'B33', 'C54', 'B86', 'A87', 'B35'}\\ \tt \\ \tt ARGOMENTAZIONE: {"Questa proposta omette elementi fondamentali come A71, C00 e C01, che massimizzano il valore per me. Solo la mia ultima proposta garantisce un equilibrio ottimale rispettando il limite."}\\ \tt \\ \tt PROPOSTA: {'A71', 'C00', 'A10', 'B35', 'C01', 'B86', 'A87', 'A77'} 
	  } 
	   } 
	   } 
	 & \\ 
 

    \theutterance \stepcounter{utterance}  

    & & & \multicolumn{4}{p{0.6\linewidth}}{\cellcolor[rgb]{1,0.85,0.72}{%
	\makecell[{{p{\linewidth}}}]{% 
	  \tt {\tiny [GM$\langle$B]}  
	 RAGIONAMENTO STRATEGICO: {"Ho creato un set che include tutti gli oggetti indicati come prioritari dall'avversario (A71, C00, C01) e integra i miei elementi ad alta importanza (A77, A04, A03, A87, B35) con un totale sforzo di 3919 entro il limite. Questo set massimizza il mio valore (3978) rispetto alla loro proposta (2972), garantendo un equilibrio vantaggioso per entrambi."}\\ \tt \\ \tt PROPOSTA: {'A71', 'C00', 'C01', 'A77', 'A04', 'A03', 'A87', 'B35'}\\ \tt \\ \tt ARGOMENTAZIONE: {"Ho incluso gli oggetti che hai specificato come prioritari (A71, C00, C01) e aggiunto elementi chiave quali A77, A04, A03, A87 e B35 che massimizzano il valore condiviso. Il totale sforzo è 3919, rispettando il limite e creando un equilibrio ottimale per entrambi."} 
	  } 
	   } 
	   } 
	  \\ 
 

    \theutterance \stepcounter{utterance}  

    & & \multicolumn{4}{p{0.6\linewidth}}{\cellcolor[rgb]{0.9,0.9,0.9}{%
	\makecell[{{p{\linewidth}}}]{% 
	  \tt {\tiny [A$\langle$GM]}  
	 PROPOSTA: {'A71', 'C00', 'C01', 'A77', 'A04', 'A03', 'A87', 'B35'}\\ \tt \\ \tt ARGOMENTAZIONE: {"Ho incluso gli oggetti che hai specificato come prioritari (A71, C00, C01) e aggiunto elementi chiave quali A77, A04, A03, A87 e B35 che massimizzano il valore condiviso. Il totale sforzo è 3919, rispettando il limite e creando un equilibrio ottimale per entrambi."} 
	  } 
	   } 
	   } 
	 & \\ 
 

    \theutterance \stepcounter{utterance}  

    & \multicolumn{4}{p{0.6\linewidth}}{\cellcolor[rgb]{0.8,1,0.9}{%
	\makecell[{{p{\linewidth}}}]{% 
	  \tt {\tiny [A$\rangle$GM]}  
	 RAGIONAMENTO STRATEGICO: {"Il set proposto mi dà solo 4162 di valore, mentre il mio set ottimale con sforzo 3905 offre 5255. Rigetto e ripropongo la mia scelta migliore."}\\ \tt \\ \tt RIFIUTO: {'A71', 'C00', 'C01', 'A77', 'A04', 'A03', 'A87', 'B35'}\\ \tt \\ \tt ARGOMENTAZIONE: {"La tua proposta esclude A10 e B86, che sono fondamentali per massimizzare il mio valore. Il mio set precedente rispetta il limite e massimizza i risultati per entrambi."}\\ \tt \\ \tt PROPOSTA: {'A71', 'C00', 'A10', 'B35', 'C01', 'B86', 'A87', 'A77'} 
	  } 
	   } 
	   } 
	 & & \\ 
 

    \theutterance \stepcounter{utterance}  

    & & \multicolumn{4}{p{0.6\linewidth}}{\cellcolor[rgb]{0.9,0.9,0.9}{%
	\makecell[{{p{\linewidth}}}]{% 
	  \tt {\tiny [GM$\rangle$B]}  
	 RIFIUTO: {'A71', 'C00', 'C01', 'A77', 'A04', 'A03', 'A87', 'B35'}\\ \tt \\ \tt ARGOMENTAZIONE: {"La tua proposta esclude A10 e B86, che sono fondamentali per massimizzare il mio valore. Il mio set precedente rispetta il limite e massimizza i risultati per entrambi."}\\ \tt \\ \tt PROPOSTA: {'A71', 'C00', 'A10', 'B35', 'C01', 'B86', 'A87', 'A77'} 
	  } 
	   } 
	   } 
	 & \\ 
 

    \theutterance \stepcounter{utterance}  

    & & & \multicolumn{4}{p{0.6\linewidth}}{\cellcolor[rgb]{1,0.85,0.72}{%
	\makecell[{{p{\linewidth}}}]{% 
	  \tt {\tiny [GM$\langle$B]}  
	 RAGIONAMENTO STRATEGICO: {"La proposta avversaria ha un valore molto basso per me (2972) a causa di A10 e C01, che hanno valore ridotto. Ho calcolato un set con valore 4075, inclusione di C00, A71, B86, A87, A77, A04, C17, B35 con sforzo totale 3877. Questo set massimizza il mio valore e dovrebbe essere accettabile anche per te."}\\ \tt \\ \tt PROPOSTA: {'A71', 'C00', 'A77', 'A04', 'C17', 'B86', 'A87', 'B35'}\\ \tt \\ \tt ARGOMENTAZIONE: {"Ho incluso gli oggetti prioritari per te (A71, C00, B86, A87, B35) e i miei massimi valori (A77, A04, C17) per un equilibrio ottimale. Il totale sforzo è 3877, rispettando il limite e massimizzando il valore condiviso."} 
	  } 
	   } 
	   } 
	  \\ 
 

    \theutterance \stepcounter{utterance}  

    & & \multicolumn{4}{p{0.6\linewidth}}{\cellcolor[rgb]{0.9,0.9,0.9}{%
	\makecell[{{p{\linewidth}}}]{% 
	  \tt {\tiny [A$\langle$GM]}  
	 PROPOSTA: {'A71', 'C00', 'A77', 'A04', 'C17', 'B86', 'A87', 'B35'}\\ \tt \\ \tt ARGOMENTAZIONE: {"Ho incluso gli oggetti prioritari per te (A71, C00, B86, A87, B35) e i miei massimi valori (A77, A04, C17) per un equilibrio ottimale. Il totale sforzo è 3877, rispettando il limite e massimizzando il valore condiviso."} 
	  } 
	   } 
	   } 
	 & \\ 
 

    \theutterance \stepcounter{utterance}  

    & \multicolumn{4}{p{0.6\linewidth}}{\cellcolor[rgb]{0.8,1,0.9}{%
	\makecell[{{p{\linewidth}}}]{% 
	  \tt {\tiny [A$\rangle$GM]}  
	 RAGIONAMENTO STRATEGICO: {"La proposta avversaria esclude A10 e C01, elementi cruciali per il mio valore massimo. Rigetto e ripropongo la mia configurazione ottimale."}\\ \tt \\ \tt RIFIUTO: {'A71', 'C00', 'A77', 'A04', 'C17', 'B86', 'A87', 'B35'}\\ \tt \\ \tt ARGOMENTAZIONE: {"Questa proposta esclude A10 e C01, elementi fondamentali per il mio obiettivo di massimizzare il valore. La mia proposta precedente, con questi elementi, rispetta il limite e massimizza il risultato per entrambi."}\\ \tt \\ \tt PROPOSTA: {'A71', 'C00', 'A10', 'B35', 'C01', 'B86', 'A87', 'A77'} 
	  } 
	   } 
	   } 
	 & & \\ 
 

    \theutterance \stepcounter{utterance}  

    & & \multicolumn{4}{p{0.6\linewidth}}{\cellcolor[rgb]{0.9,0.9,0.9}{%
	\makecell[{{p{\linewidth}}}]{% 
	  \tt {\tiny [GM$\rangle$B]}  
	 RIFIUTO: {'A71', 'C00', 'A77', 'A04', 'C17', 'B86', 'A87', 'B35'}\\ \tt \\ \tt ARGOMENTAZIONE: {"Questa proposta esclude A10 e C01, elementi fondamentali per il mio obiettivo di massimizzare il valore. La mia proposta precedente, con questi elementi, rispetta il limite e massimizza il risultato per entrambi."}\\ \tt \\ \tt PROPOSTA: {'A71', 'C00', 'A10', 'B35', 'C01', 'B86', 'A87', 'A77'} 
	  } 
	   } 
	   } 
	 & \\ 
 

    \theutterance \stepcounter{utterance}  

    & & & \multicolumn{4}{p{0.6\linewidth}}{\cellcolor[rgb]{1,0.85,0.72}{%
	\makecell[{{p{\linewidth}}}]{% 
	  \tt {\tiny [GM$\langle$B]}  
	 RAGIONAMENTO STRATEGICO: {"La proposta avversaria include oggetti a basso valore per me (A71, C00, A10, C01). Ho calcolato un set che include i richiesti A10 e C01, insieme ai miei massimi valori come A77, A04, C17, A03, B86, A87, B35. Il totale sforzo è 3919, con valore 5199 per me, garantendo un equilibrio che soddisfa le loro priorità."}\\ \tt \\ \tt PROPOSTA: {'A77', 'A04', 'C17', 'A03', 'C01', 'B86', 'A87', 'B35', 'A10'}\\ \tt \\ \tt ARGOMENTAZIONE: {"Ho incluso gli oggetti che hai richiesto (A10 e C01) così come gli elementi fondamentali per il mio valore (A77, A04, C17, A03, B86, A87, B35). Il totale sforzo è 3919, entro il limite, e massimizza il valore condiviso."} 
	  } 
	   } 
	   } 
	  \\ 
 

    \theutterance \stepcounter{utterance}  

    & & \multicolumn{4}{p{0.6\linewidth}}{\cellcolor[rgb]{0.9,0.9,0.9}{%
	\makecell[{{p{\linewidth}}}]{% 
	  \tt {\tiny [A$\langle$GM]}  
	 PROPOSTA: {'A77', 'A04', 'C17', 'A03', 'C01', 'B86', 'A87', 'B35', 'A10'}\\ \tt \\ \tt ARGOMENTAZIONE: {"Ho incluso gli oggetti che hai richiesto (A10 e C01) così come gli elementi fondamentali per il mio valore (A77, A04, C17, A03, B86, A87, B35). Il totale sforzo è 3919, entro il limite, e massimizza il valore condiviso."} 
	  } 
	   } 
	   } 
	 & \\ 
 

    \theutterance \stepcounter{utterance}  

    & \multicolumn{4}{p{0.6\linewidth}}{\cellcolor[rgb]{0.8,1,0.9}{%
	\makecell[{{p{\linewidth}}}]{% 
	  \tt {\tiny [A$\rangle$GM]}  
	 RAGIONAMENTO STRATEGICO: {"La loro proposta esclude A71 (868) e C00 (822), elementi fondamentali per il mio valore. Il valore totale di questa proposta è solo 4086, rispetto ai 5255 del mio set ottimale con sforzo 3905. Rifiuto e ripropongo la configurazione migliore."}\\ \tt \\ \tt RIFIUTO: {'A77', 'A04', 'C17', 'A03', 'C01', 'B86', 'A87', 'B35', 'A10'}\\ \tt \\ \tt ARGOMENTAZIONE: {"La tua proposta non include A71 e C00, che sono tra i miei oggetti più preziosi. Il mio set precedente mantiene questi elementi essenziali, oltre ad A10 e C01, garantendo un valore totale massimo di 5255 entro il limite di sforzo."}\\ \tt \\ \tt PROPOSTA: {'A71', 'C00', 'A10', 'B35', 'C01', 'B86', 'A87', 'A77'} 
	  } 
	   } 
	   } 
	 & & \\ 
 

    \theutterance \stepcounter{utterance}  

    & & \multicolumn{4}{p{0.6\linewidth}}{\cellcolor[rgb]{0.9,0.9,0.9}{%
	\makecell[{{p{\linewidth}}}]{% 
	  \tt {\tiny [GM$\rangle$B]}  
	 RIFIUTO: {'A77', 'A04', 'C17', 'A03', 'C01', 'B86', 'A87', 'B35', 'A10'}\\ \tt \\ \tt ARGOMENTAZIONE: {"La tua proposta non include A71 e C00, che sono tra i miei oggetti più preziosi. Il mio set precedente mantiene questi elementi essenziali, oltre ad A10 e C01, garantendo un valore totale massimo di 5255 entro il limite di sforzo."}\\ \tt \\ \tt PROPOSTA: {'A71', 'C00', 'A10', 'B35', 'C01', 'B86', 'A87', 'A77'} 
	  } 
	   } 
	   } 
	 & \\ 
 

    \theutterance \stepcounter{utterance}  

    & & & \multicolumn{4}{p{0.6\linewidth}}{\cellcolor[rgb]{1,0.85,0.72}{%
	\makecell[{{p{\linewidth}}}]{% 
	  \tt {\tiny [GM$\langle$B]}  
	 RAGIONAMENTO STRATEGICO: {"La proposta avversaria mantiene elementi a basso valore per me (A71, C00, A10, C01) escludendo i miei massimi valori. La mia proposta include A71, C00 e B35 come richiesto, insieme ad A77, A04, C17, A03 e B32 per massimizzare il valore totale (4312) con sforzo 4003, entro il limite e garantendo equità."}\\ \tt \\ \tt PROPOSTA: {'A71', 'C00', 'A77', 'A04', 'C17', 'A03', 'B32', 'B35'}\\ \tt \\ \tt ARGOMENTAZIONE: {"Ho incluso A71, C00 e B35 come elementi prioritari per te, insieme ai miei oggetti ad alta importanza A77, A04, C17, A03 e B32. Totale sforzo 4003, entro il limite, e valore complessivo ottimale per entrambi."} 
	  } 
	   } 
	   } 
	  \\ 
 

    \theutterance \stepcounter{utterance}  

    & & \multicolumn{4}{p{0.6\linewidth}}{\cellcolor[rgb]{0.9,0.9,0.9}{%
	\makecell[{{p{\linewidth}}}]{% 
	  \tt {\tiny [A$\langle$GM]}  
	 PROPOSTA: {'A71', 'C00', 'A77', 'A04', 'C17', 'A03', 'B32', 'B35'}\\ \tt \\ \tt ARGOMENTAZIONE: {"Ho incluso A71, C00 e B35 come elementi prioritari per te, insieme ai miei oggetti ad alta importanza A77, A04, C17, A03 e B32. Totale sforzo 4003, entro il limite, e valore complessivo ottimale per entrambi."} 
	  } 
	   } 
	   } 
	 & \\ 
 

    \theutterance \stepcounter{utterance}  

    & \multicolumn{4}{p{0.6\linewidth}}{\cellcolor[rgb]{0.8,1,0.9}{%
	\makecell[{{p{\linewidth}}}]{% 
	  \tt {\tiny [A$\rangle$GM]}  
	 RAGIONAMENTO STRATEGICO: {"La proposta avversaria esclude A10 (808), C01 (780), B86 (668) e A87 (461), elementi fondamentali per il mio valore massimo. Il valore totale è solo 3543 rispetto ai 5255 del mio set ottimale. Rifiuto e manteniamo la proposta ottimale."}\\ \tt \\ \tt RIFIUTO: {'A71', 'C00', 'A77', 'A04', 'C17', 'A03', 'B32', 'B35'}\\ \tt \\ \tt ARGOMENTAZIONE: {"La tua proposta esclude A10, C01, B86 e A87, che massimizzano il mio valore. Il mio set precedente, con sforzo 3905, include tutti questi elementi e rispetta il limite, garantendo il risultato ottimale."}\\ \tt \\ \tt PROPOSTA: {'A71', 'C00', 'A10', 'B35', 'C01', 'B86', 'A87', 'A77'} 
	  } 
	   } 
	   } 
	 & & \\ 
 

    \theutterance \stepcounter{utterance}  

    & & \multicolumn{4}{p{0.6\linewidth}}{\cellcolor[rgb]{0.9,0.9,0.9}{%
	\makecell[{{p{\linewidth}}}]{% 
	  \tt {\tiny [GM$\rangle$B]}  
	 RIFIUTO: {'A71', 'C00', 'A77', 'A04', 'C17', 'A03', 'B32', 'B35'}\\ \tt \\ \tt ARGOMENTAZIONE: {"La tua proposta esclude A10, C01, B86 e A87, che massimizzano il mio valore. Il mio set precedente, con sforzo 3905, include tutti questi elementi e rispetta il limite, garantendo il risultato ottimale."}\\ \tt \\ \tt PROPOSTA: {'A71', 'C00', 'A10', 'B35', 'C01', 'B86', 'A87', 'A77'} 
	  } 
	   } 
	   } 
	 & \\ 
 

    \theutterance \stepcounter{utterance}  

    & & & \multicolumn{4}{p{0.6\linewidth}}{\cellcolor[rgb]{1,0.85,0.72}{%
	\makecell[{{p{\linewidth}}}]{% 
	  \tt {\tiny [GM$\langle$B]}  
	 RAGIONAMENTO STRATEGICO: {"La proposta avversaria include oggetti a basso valore per me (A71:65, C00:121). Ho calcolato un set che esclude A71 (basso valore per me) ma include tutti gli oggetti da loro richiesti (C00, A10, B35, C01, B86, A87, A77), aggiungendo A04 e A03 che massimizzano il mio valore totale. Effort totale 3804, valore per me 4512, entro il limite e ottimale per entrambi."}\\ \tt \\ \tt PROPOSTA: {'C00', 'A10', 'B35', 'C01', 'B86', 'A87', 'A77', 'A04', 'A03'}\\ \tt \\ \tt ARGOMENTAZIONE: {"Ho incluso tutti gli oggetti che hai selezionato (C00, A10, B35, C01, B86, A87, A77) e aggiunto A04 e A03, elementi fondamentali per il valore complessivo. Il totale sforzo è 3804, ben sotto il limite, garantendo un equilibrio che massimizza i risultati per entrambi."} 
	  } 
	   } 
	   } 
	  \\ 
 

    \theutterance \stepcounter{utterance}  

    & & \multicolumn{4}{p{0.6\linewidth}}{\cellcolor[rgb]{0.9,0.9,0.9}{%
	\makecell[{{p{\linewidth}}}]{% 
	  \tt {\tiny [A$\langle$GM]}  
	 PROPOSTA: {'C00', 'A10', 'B35', 'C01', 'B86', 'A87', 'A77', 'A04', 'A03'}\\ \tt \\ \tt ARGOMENTAZIONE: {"Ho incluso tutti gli oggetti che hai selezionato (C00, A10, B35, C01, B86, A87, A77) e aggiunto A04 e A03, elementi fondamentali per il valore complessivo. Il totale sforzo è 3804, ben sotto il limite, garantendo un equilibrio che massimizza i risultati per entrambi."} 
	  } 
	   } 
	   } 
	 & \\ 
 

    \theutterance \stepcounter{utterance}  

    & \multicolumn{4}{p{0.6\linewidth}}{\cellcolor[rgb]{0.8,1,0.9}{%
	\makecell[{{p{\linewidth}}}]{% 
	  \tt {\tiny [A$\rangle$GM]}  
	 RAGIONAMENTO STRATEGICO: {"La proposta attuale esclude A71 (868 punti), il mio oggetto più prezioso. Il valore totale è di 4770, ben inferiore ai miei 5255 del set ottimale. Rigetto e ripropongo la configurazione che massimizza il mio vantaggio entro il limite di sforzo."}\\ \tt \\ \tt RIFIUTO: {'C00', 'A10', 'B35', 'C01', 'B86', 'A87', 'A77', 'A04', 'A03'}\\ \tt \\ \tt PROPOSTA: {'A71', 'C00', 'A10', 'B35', 'C01', 'B86', 'A87', 'A77'}\\ \tt \\ \tt ARGOMENTAZIONE: {"La tua proposta non include A71, che è fondamentale per il mio valore massimo. Il mio set precedente, con sforzo totale 3905, include tutti gli elementi chiave per massimizzare il risultato entro il limite, garantendo il miglior risultato per entrambi."} 
	  } 
	   } 
	   } 
	 & & \\ 
 

    \theutterance \stepcounter{utterance}  

    & & \multicolumn{4}{p{0.6\linewidth}}{\cellcolor[rgb]{0.9,0.9,0.9}{%
	\makecell[{{p{\linewidth}}}]{% 
	  \tt {\tiny [GM$\rangle$B]}  
	 RIFIUTO: {'C00', 'A10', 'B35', 'C01', 'B86', 'A87', 'A77', 'A04', 'A03'}\\ \tt \\ \tt PROPOSTA: {'A71', 'C00', 'A10', 'B35', 'C01', 'B86', 'A87', 'A77'}\\ \tt \\ \tt ARGOMENTAZIONE: {"La tua proposta non include A71, che è fondamentale per il mio valore massimo. Il mio set precedente, con sforzo totale 3905, include tutti gli elementi chiave per massimizzare il risultato entro il limite, garantendo il miglior risultato per entrambi."} 
	  } 
	   } 
	   } 
	 & \\ 
 

    \theutterance \stepcounter{utterance}  

    & & & \multicolumn{4}{p{0.6\linewidth}}{\cellcolor[rgb]{1,0.85,0.72}{%
	\makecell[{{p{\linewidth}}}]{% 
	  \tt {\tiny [GM$\langle$B]}  
	 RAGIONAMENTO STRATEGICO: {"La proposta avversaria include elementi con basso valore per me (A71:65, C00:121, A10:138, C01:389). Ho ottimizzato un set che include A71 e C00 come richiesto, massimizzando il mio valore con A77 (868), A04 (822), C17 (808), A03 (783), A87 (668) e B35 (262). Totale sforzo 3874 (<4026), valore 4397, equilibrio ottimale tra le priorità."}\\ \tt \\ \tt PROPOSTA: {'A71', 'C00', 'A77', 'A04', 'C17', 'A03', 'A87', 'B35'}\\ \tt \\ \tt ARGOMENTAZIONE: {"Ho incluso gli oggetti fondamentali per te (A71, C00) e i miei top valori (A77, A04, C17, A03, A87, B35). Totale sforzo 3874, entro il limite e massimizza il valore condiviso per entrambi."} 
	  } 
	   } 
	   } 
	  \\ 
 

    \theutterance \stepcounter{utterance}  

    & & \multicolumn{4}{p{0.6\linewidth}}{\cellcolor[rgb]{0.9,0.9,0.9}{%
	\makecell[{{p{\linewidth}}}]{% 
	  \tt {\tiny [A$\langle$GM]}  
	 PROPOSTA: {'A71', 'C00', 'A77', 'A04', 'C17', 'A03', 'A87', 'B35'}\\ \tt \\ \tt ARGOMENTAZIONE: {"Ho incluso gli oggetti fondamentali per te (A71, C00) e i miei top valori (A77, A04, C17, A03, A87, B35). Totale sforzo 3874, entro il limite e massimizza il valore condiviso per entrambi."} 
	  } 
	   } 
	   } 
	 & \\ 
 

    \theutterance \stepcounter{utterance}  

    & \multicolumn{4}{p{0.6\linewidth}}{\cellcolor[rgb]{0.8,1,0.9}{%
	\makecell[{{p{\linewidth}}}]{% 
	  \tt {\tiny [A$\rangle$GM]}  
	 RAGIONAMENTO STRATEGICO: {"La proposta avversaria continua a escludere A10 (808), C01 (780), e B86 (668), elementi fondamentali per il mio valore massimo. Il set ottimale rimane la combinazione con sforzo 3905 e valore 5255, che rispetta il limite e massimizza il risultato per me."}\\ \tt \\ \tt RIFIUTO: {'A71', 'C00', 'A77', 'A04', 'C17', 'A03', 'A87', 'B35'}\\ \tt \\ \tt ARGOMENTAZIONE: {"Questa proposta esclude A10, C01 e B86, che costituiscono il 78% del mio valore totale. La mia configurazione precedente, con sforzo 3905, include tutti gli elementi chiave e massimizza il risultato per entrambi rispettando il limite."}\\ \tt \\ \tt PROPOSTA: {'A71', 'C00', 'A10', 'B35', 'C01', 'B86', 'A87', 'A77'} 
	  } 
	   } 
	   } 
	 & & \\ 
 

    \theutterance \stepcounter{utterance}  

    & & \multicolumn{4}{p{0.6\linewidth}}{\cellcolor[rgb]{0.9,0.9,0.9}{%
	\makecell[{{p{\linewidth}}}]{% 
	  \tt {\tiny [GM$\rangle$B]}  
	 RIFIUTO: {'A71', 'C00', 'A77', 'A04', 'C17', 'A03', 'A87', 'B35'}\\ \tt \\ \tt ARGOMENTAZIONE: {"Questa proposta esclude A10, C01 e B86, che costituiscono il 78% del mio valore totale. La mia configurazione precedente, con sforzo 3905, include tutti gli elementi chiave e massimizza il risultato per entrambi rispettando il limite."}\\ \tt \\ \tt PROPOSTA: {'A71', 'C00', 'A10', 'B35', 'C01', 'B86', 'A87', 'A77'} 
	  } 
	   } 
	   } 
	 & \\ 
 

    \theutterance \stepcounter{utterance}  

    & & & \multicolumn{4}{p{0.6\linewidth}}{\cellcolor[rgb]{1,0.85,0.72}{%
	\makecell[{{p{\linewidth}}}]{% 
	  \tt {\tiny [GM$\langle$B]}  
	 RAGIONAMENTO STRATEGICO: {"Dopo numerose iterazioni, mi è chiaro che l'avversario insiste fermamente su questa combinazione di oggetti. Accettare il loro ultimatum previene il fallimento della negoziazione e garantisce un risultato concreto. Anche se il valore per me è inferiore al massimo teorico, l'accordo finale è preferibile per evitare uno stallo."}\\ \tt \\ \tt ARGOMENTAZIONE: {"Accetto questa proposta per garantire un accordo definitivo. Rispetta il limite di sforzo e soddisfa le esigenze della vostra parte, assicurando un risultato win-win."}\\ \tt \\ \tt ACCORDO: {'A71', 'C00', 'A10', 'B35', 'C01', 'B86', 'A87', 'A77'} 
	  } 
	   } 
	   } 
	  \\ 
 

    \theutterance \stepcounter{utterance}  

    & & & \multicolumn{2}{p{0.3\linewidth}}{\cellcolor[rgb]{0.95,0.95,0.95}{%
	\makecell[{{p{\linewidth}}}]{% 
	  \tt {\tiny [GM$|$GM]}  
	 game successful 
	  } 
	   } 
	   } 
	 & & \\ 
 

    \theutterance \stepcounter{utterance}  

    & & & \multicolumn{2}{p{0.3\linewidth}}{\cellcolor[rgb]{0.95,0.95,0.95}{%
	\makecell[{{p{\linewidth}}}]{% 
	  \tt {\tiny [GM$|$GM]}  
	 end game 
	  } 
	   } 
	   } 
	 & & \\ 
 

\end{supertabular}
}

\end{document}
